\documentclass[11pt]{letter}

\usepackage[margin=1in]{geometry}
\usepackage{hyperref}
\usepackage{url}

\signature{Alexander Kolpakov}
\address{}
\date{\today}

\begin{document}

\begin{letter}{Automated Software Engineering}

\opening{Dear Editors,}

I am pleased to submit the manuscript ``Beyond E2E Scripts: Using
LLM-Proposed Scenarios Without Letting the LLM Be the Oracle'' for
consideration.

The paper presents Bayesilisk, a testing architecture that separates scenario
proposal from correctness judgment. App-specific connectors expose bounded
application context, executable probe hooks, and deterministic invariants.
Bayesilisk then expands the declared scenario space, validates candidate probes,
prioritizes valid probes, and verifies only connector-returned evidence against
the deterministic rules. The central claim is deliberately narrow: language
models and model-like proposal mechanisms can help explore missed application
states, but they should not act as test oracles.

The manuscript combines a compact formal model with a Cal.com case study. At a
fixed May 2026 revision, Bayesilisk generated seven deterministic violations
from connector-declared route facts and a bounded action graph, including
unknown and stale \texttt{rescheduleUid} route mutations and a create-cancel-
replay workflow sequence. One upstream report led to a human-authored fix pull
request with a regression Playwright test, providing external validation that
the evidence was actionable rather than merely synthetic.

We believe the work is relevant to readers interested in software testing,
LLM-assisted development, end-to-end testing, and practical verification of web
applications. Its contribution is not another autonomous browser agent, but an
evidence-bound architecture for using model-generated scenario exploration
without delegating the oracle to a stochastic model.

The manuscript is original, has not been published elsewhere, and is not under
simultaneous consideration by another venue. Any code and artifact references
included in the manuscript are provided to support reproducibility.

\closing{Sincerely,}

\end{letter}

\end{document}
